\chapter{Analysis and Interpretability}
\label{ch:analysis}

\section{Overview}

This chapter interprets the Phase~1 and Phase~2 results reported in
\Cref{ch:results}. The analysis focuses on three patterns that emerged from
scaled experiments on Metamon Gen~9 OU human replays
\begin{enumerate}
  \item move-belief accuracy improves with data scale while policy metrics plateau
    across architectures,
  \item opponent-team generalization gaps shrink with corpus size for every
    model, and
  \item turn-only behavioral cloning matches or exceeds matchup- and
    belief-augmented variants on action prediction.
\end{enumerate}

\section{What the Thesis Shows Versus What It Does Not}

The offline metrics support a scientific claim about learning, not a product
claim about play. Nontrivial move belief and far-above-chance action match show
that the representation and training pipeline recover structure from human logs.
They do not imply that the agent would win live battles against real players.
A model that matches the human action on roughly one turn in five still errs on
most turns, and a single bad switch often decides a game. Interpreting the
thesis as a failed competitive bot would misread the research questions.

The negative result on belief-to-policy transfer is equally important. Explicit
belief heads improve decodable move prediction without improving action match.
That finding warns against assuming that auxiliary opponent modeling will
automatically improve imitation policies.

\section{Scaling Belief Versus Scaling Policy}

The \texttt{ActiveBeliefTransformer} increased validation move-slot accuracy from
15.5\% on a 500-battle pilot export to 28.1\% on a 2\,000-battle export, while
validation action accuracy rose from 8.5\% to 20.2\% over the same period. The
belief head therefore benefits directly from additional labeled trajectories.

Action accuracy for all three Phase~2 models converged near 20\% on random battle
holdout (Split~A) after scaling, with turn-only BC slightly ahead. Auxiliary
belief losses add parameters and optimization pressure but did not produce a
policy advantage at $d=128$ and twenty epochs. Two explanations are consistent
with the data.

\textbf{Shared trunk sufficiency.} The hierarchical turn encoder may already
compress history into an embedding useful for both action and hidden-set
prediction. Explicit belief heads learn to \emph{decode} that embedding for
supervised targets without changing the trunk enough to improve actions.

\textbf{Task interference.} Multi-task training divides gradient signal among
move, item, ability, and tera heads. If $\lambda_\mathrm{bel}$ is not carefully
tuned, belief terms may slow action learning without improving generalization.
Belief fusion into the policy input was tested and did not help
(\Cref{ch:results}), so the remaining open knob is a careful $\lambda_\mathrm{bel}$
sweep rather than fusion alone.

\section{Why the Team Generalization Gap Collapsed}

On the 500-battle corpus, Split~B showed a 1.8--1.9 point gap between ID train
teams and OOD teams. At 2\,000 battles the gap fell below half a point for every
architecture, including turn-only BC.

This pattern suggests that the pilot gap was driven largely by \emph{data
coverage} rather than architectural failure. With only 500 battles, many
held-out team keys had few or no nearby examples in embedding space. Expanding
to 2\,000 battles increased overlap in species co-occurrence patterns and move
distributions, so OOD team turns became easier to match even without explicit
matchup modules.

Matchup BC did not outperform turn-only BC on Split~B at scale. The
\texttt{TeamMatchupEncoder} pools six species embeddings per side and adds them
to the final turn vector. Much of that information is already present in the
opponent party tokens within each turn's thirteen-token block. The explicit
module may therefore be redundant given current tokenization.

Split~B accuracy ($\approx$42--45\%) exceeds Split~A accuracy ($\approx$20\%)
because the metrics condition on different splits. Split~A withholds entire
battles. Split~B groups turns by roster key while many action contexts remain
familiar. Reporting both is essential to avoid over-interpreting high team-split
numbers as solved imitation.

\section{Linking Beliefs to Actions}

The original plan correlated per-battle belief accuracy with per-battle action
match on held-out teams. Species belief was unavailable in Gen~9 OU because team
preview reveals all opponent species. Move-belief correlation analysis is still
informative.

At scaled training, \texttt{active\_belief} reached 28.1\% move belief and 44.4\%
OOD team action match, while turn-only BC reached 45.1\% OOD action match with no
belief head. Higher move-belief accuracy therefore did not coincide with better
OOD policy performance in this run.

Linear probes on frozen \texttt{human\_bc\_belief} trunk embeddings tell a
sharper story (\Cref{tab:trunk-probe}, script
\texttt{scripts/probe\_trunk\_belief.py}). A fifteen-epoch linear probe on the
final turn embedding reached only \textbf{10.5\%} validation move-slot accuracy,
versus \textbf{27.7\%} for \texttt{active\_belief} move heads on the same split.
The BC trunk therefore does \emph{not} expose hidden moves through an easy linear
readout. Explicit belief heads and multi-task training learn representations that
decode move labels far better than post-hoc probing.

Action accuracy nevertheless remains similar across models. Belief-specific
features appear to live in a subspace that the shared action head does not
exploit, or multi-task training reshapes the trunk in ways that help belief
decoding without improving the action logit channel used at inference.

\begin{table}[ht]
  \centering
  \caption{Move-slot belief accuracy on validation battles (same split).}
  \label{tab:trunk-probe}
  \begin{tabular}{lc}
    \toprule
    Method & Val.\ move belief \\
    \midrule
    Linear probe on frozen BC trunk & 10.5\% \\
    ActiveBelief move heads & 27.7\% \\
    \bottomrule
  \end{tabular}
\end{table}

\section{Implicit Belief in the Turn-Only Trunk}

Turn-only BC and ActiveBelief achieved nearly identical OOD team accuracy
(45.1\% versus 44.4\%) despite a large difference in supervised move-belief
accuracy. Linear probes on the frozen BC trunk reached only 10.5\% move belief
(\Cref{tab:trunk-probe}), so implicit linear encodings of unrevealed moves are
weak. Explicit belief heads add substantial decodable structure without
translating into better action match at current scale.

The thirteen-token layout exposes opponent party state, field conditions, and
revealed moves each turn. Causal attention over eight turns integrates switch
patterns, damage clues, and item triggers. That design aligns with the thesis
motivation for structured partial observability without a separate belief filter.

\section{Gen~9 OU and Belief Target Selection}

Team preview fundamentally changed the belief task relative to older formats
where species were hidden. The implementation pivot to active-opponent move,
item, ability, and tera labels was necessary for non-empty supervision.

Species-centric hypotheses (H1 as originally framed) do not apply verbatim to this
benchmark slice. Move-belief results should be read as evidence that history-
conditioned transformers and classical baselines infer unrevealed moves, not that
the full hidden-state posterior is solved.

\section{Interpretability Directions Not Executed}

Shallow decision trees on hand-crafted turn features and attention visualizations
were planned to compare interpretable clues against neural errors. Those studies
were not completed for this draft. Expected error modes remain relevant for
future case studies
\begin{itemize}
  \item rare move sets on held-out teams,
  \item ambiguous damage ranges before stat clues accumulate,
  \item deliberate unconventional sets underrepresented in training,
  \item late-game turns with truncated history relative to the eight-turn window.
\end{itemize}

The quantitative baselines already provide the main interpretability result for
RQ1. A latest-turn MLP beating the transformer suggests that most move-belief
signal lives in the current observation rather than in long-horizon attention
patterns. Attention case studies would refine that story but are not required to
support the claim.

\section{Takeaways for RQ1--RQ4}

\textbf{RQ1.} Transformers beat GRU on move belief but not a latest-turn MLP,
which reaches 35.5\% validation accuracy versus 28.1\% for
\texttt{active\_belief}.

\textbf{RQ2.} Belief augmentation did not improve action match over turn-only BC
at current scale. Offline imitation success does not imply live competitiveness
against humans.

\textbf{RQ3.} Combinatorial team shift was mitigated primarily by scaling data,
not by matchup or belief modules.

\textbf{RQ4.} Hierarchical turn structure appears sufficient for strong BC at this
scale. Belief fusion into the policy head did not improve action match over
parallel heads alone (19.6\% versus 20.2\% Split~A). Flat-sequence and
context-window sweeps remain open.

These interpretations motivate the limitations and future directions in
\Cref{ch:limitations} and the revised conclusions in \Cref{ch:conclusion}.
